\section{Lecture 19 (February 28, 2025, est.): The Frobenius functor}
\label{lec19}
\noindent\textit{Scribe: Guillermo Sanmarco}

Throughout this lecture, $k$ is algebraically closed of characteristic $p>0$,
and $\mathcal C$ is a semisimple symmetric tensor category. We write
$\overline{\mathcal A}$ for the semisimplification of a symmetric category
$\mathcal A$ by negligible morphisms. For the cyclic group $C_p$, let $J_r$
be the indecomposable representation of dimension $r$, and let $L_r$ be its
image in $\mathsf{Ver}_p$. Thus $L_p=0$ and $L_1=\mathbf1$.

\subsection{Equivariant objects and semisimplification}
If a finite group $G$ acts trivially on $\mathcal C$, an object of
$\mathcal C^G$ is a pair $(X,\rho)$ with $\rho:G\to\Aut_{\mathcal C}(X)$.
The morphisms commute with the specified actions. There is a symmetric
equivalence
\[
 \mathcal C^G\simeq\mathcal C\boxtimes\Rep_k(G).
\]
Indeed, write $X=\bigoplus_S S\otimes_k M_S$, where $S$ runs through the
simple objects of $\mathcal C$ and $M_S=\Hom(S,X)$. Since
$\End(S)=k$ and there are no morphisms between distinct simples, the
$G$-action is precisely a representation on each multiplicity space $M_S$.

\begin{proposition}\label{lec19-equivariant-semisimplification}
There is a symmetric equivalence
\[
 \overline{\mathcal C^G}
 \simeq\mathcal C\boxtimes\overline{\Rep_k(G)}.
\]
\end{proposition}
\begin{proof}
The indecomposables of $\mathcal C^G$ are $S\boxtimes M$, with $S$ simple
and $M$ an \emph{indecomposable} $G$-module. Their categorical dimensions are
$\dim(S)\dim_k(M)$. A simple object in a semisimple rigid category has
nonzero categorical dimension: the two one-dimensional spaces
$\Hom(\mathbf1,S\otimes S^*)$ and $\Hom(S\otimes S^*,\mathbf1)$ pair
nondegenerately, because the unit occurs with multiplicity one.
Consequently $S\boxtimes M$ has dimension zero exactly when $M$ does.
The trace pairing on the $S$-isotypic block is the ordinary trace pairing
on the multiplicity spaces multiplied by the nonzero scalar $\dim(S)$.
Thus taking its radical, and then the quotient, acts separately on each
$G$-module block. This gives the claimed equivalence, compatibly with tensor
products and symmetry.
\end{proof}

\begin{remark}
It is simple objects, rather than all nonzero objects, that have nonzero
dimension here. For example, $\mathbf1^{\oplus p}$ is nonzero but has
dimension zero in characteristic $p$.
\end{remark}

\subsection{Tensor powers and additivity}
The symmetry defines an action of $S_p$ on $X^{\otimes p}$. Restrict it to
$C_p=\langle\sigma\rangle$, where $\sigma=(1\,2\,\cdots\,p)$. This gives a
functor
\[
 P:\mathcal C\longrightarrow\mathcal C^{C_p},\qquad
 X\longmapsto(X^{\otimes p},\sigma),\quad f\longmapsto f^{\otimes p}.
\]
The canonical regrouping of tensor factors makes $P$ symmetric monoidal,
but $P$ is not additive. Compose it with semisimplification to obtain
\[
 \mathrm{Fr}_0:\mathcal C\longrightarrow
 \overline{\mathcal C^{C_p}}\simeq
 \mathcal C\boxtimes\mathsf{Ver}_p.
\]

\begin{lemma}\label{lec19-additivity}
The functor $\mathrm{Fr}_0$ is additive.
\end{lemma}
\begin{proof}
For $f,g:X\to Y$, expand $(f+g)^{\otimes p}$. Besides $f^{\otimes p}$ and
$g^{\otimes p}$, every term is a word in $f$ and $g$ containing both letters.
The cyclic group acts freely on these words, since $p$ is prime. The sum
over one orbit has the form
\[
 T=\sum_{i=0}^{p-1}\sigma_Y^i h\sigma_X^{-i},
 \qquad h=h_1\otimes\cdots\otimes h_p.
\]
It is $C_p$-equivariant. For an equivariant map
$u:Y^{\otimes p}\to X^{\otimes p}$, naturality and cyclicity of the trace give
\[
 \Tr(\sigma_Y^i h\sigma_X^{-i}u)
 =\Tr(h\sigma_X^{-i}u\sigma_Y^i)=\Tr(hu).
\]
Hence $\Tr(Tu)=p\Tr(hu)=0$. Each orbit sum is negligible, so the mixed terms
vanish after semisimplification. This proves additivity on morphisms.
Since $P(0)=0$, preservation of biproducts follows as well.
\end{proof}

\subsection{The scalar twist}
Although additive, $\mathrm{Fr}_0$ satisfies
$\mathrm{Fr}_0(\lambda f)=\lambda^p\mathrm{Fr}_0(f)$. Define
$\mathcal C^{(1)}$ to have the same underlying additive symmetric monoidal
category as $\mathcal C$, with scalar multiplication
\[
 \lambda\mathbin{\cdot_{(1)}} f=\lambda^p f.
\]
Because $k$ is perfect, this is again a tensor category over $k$.
The category $\mathsf{Ver}_p$ is defined over $\mathbb F_p$, and therefore
$\mathsf{Ver}_p^{(1)}\simeq\mathsf{Ver}_p$.

\begin{definition}
The \emph{Frobenius functor} is the resulting $k$-linear symmetric monoidal
functor
\[
 \mathrm{Fr}:\mathcal C\longrightarrow
 \mathcal C^{(1)}\boxtimes\mathsf{Ver}_p.
\]
\end{definition}
It is exact because every exact sequence in $\mathcal C$ splits. This
argument uses semisimplicity; exactness will be a genuine condition in the
last lecture. The construction and additivity argument are due to Ostrik
\cite[Section 3.3]{OstrikVerlinde}.

\subsection{The example $\mathsf{Ver}_5^+$}
Let $X=L_3$ in $\mathsf{Ver}_5^+$. Its fusion rule is
$X^{\otimes2}=\mathbf1\oplus X$, so
\[
 X^{\otimes5}=\mathbf1^{\oplus3}\oplus X^{\oplus5}.
\]
As an equivariant object this has the form
$\mathbf1\boxtimes M_1\oplus X\boxtimes M_X$, with
$\dim_k M_1=3$ and $\dim_k M_X=5$. Put
$\varphi=(1+\sqrt5)/2=\operatorname{FPdim}(X)$. Tensor functors between
fusion categories preserve Frobenius--Perron dimensions, hence
\[
 \varphi=\operatorname{FPdim}(\overline{M_1})
       +\varphi\operatorname{FPdim}(\overline{M_X}).
\]
The possible decompositions of a three-dimensional $C_5$-module are
$J_3$, $J_2\oplus J_1$, and $J_1^{\oplus3}$. Their semisimplifications have
Frobenius--Perron dimensions $\varphi$, $\varphi+1$, and $3$, respectively.
Thus equality forces $M_1=J_3$ and $\overline{M_X}=0$. Since $M_X$ has
dimension five, it must be $J_5$. Therefore
\[
 \mathrm{Fr}(X)=\mathbf1\boxtimes L_3.
\]

\subsection{The induced operation on the Grothendieck ring}
Write $K_0(\mathcal C)$ for the Grothendieck ring; a Frobenius twist does
not change this ring. Categorical dimension defines a ring map
\[
 \phi:K_0(\mathcal C)\otimes_{\mathbb Z}K_0(\mathsf{Ver}_p)
 \longrightarrow K_0(\mathcal C)\otimes_{\mathbb Z}k,
 \qquad [Y]\otimes[L_r]\longmapsto[Y]\otimes r.
\]
\begin{lemma}\label{lec19-grothendieck-frobenius}
For every $X\in\mathcal C$, one has $\phi([\mathrm{Fr}(X)])=[X]^p$.
\end{lemma}
\begin{proof}
Decompose $P(X)$ as $\bigoplus_S S\boxtimes M_S$. Forgetting the cyclic
action gives $[X]^p=\sum_S\dim_k(M_S)[S]$. On semisimplifying, each
$J_r$ with $r<p$ contributes categorical dimension $r$, while $J_p$ is
killed and its dimension is already zero in $k$. The two expressions agree.
\end{proof}
